\GAGASection{\S 2. The analytic space associated with an algebraic variety}{gaga:section:2}

In what follows, we shall consider algebraic varieties over the field $\mathbf{C}$. Such a variety will be endowed with two topologies: the ``usual'' topology and the Zariski topology. To avoid confusion, we shall prefix notions relating to the latter topology with the letter Z; for example, ``Z-open'' will mean ``open for the Zariski topology.''

\GAGASubsection{5. Definition of the analytic space associated with an algebraic variety}{gaga:no:5}

The possibility of endowing every algebraic variety with a structure of analytic space follows from the next lemma.

\GAGATarget{gaga:lemma:1}\textsc{Lemma 1}. --- \GAGATarget{gaga:lemma:1:item:a}a) \emph{The Zariski topology on $\mathbf{C}^n$ is coarser than the usual topology.}

\GAGATarget{gaga:lemma:1:item:b}b) \emph{Every Z-locally closed subset of $\mathbf{C}^n$ is analytic.}

\GAGATarget{gaga:lemma:1:item:c}c) \emph{If $\mathrm{U}$ and $\mathrm{U}'$ are Z-locally closed subsets of $\mathbf{C}^n$ and $\mathbf{C}^{n'}$, respectively, and if $f:\mathrm{U}\to\mathrm{U}'$ is a regular map, then $f$ is holomorphic.}

\GAGATarget{gaga:lemma:1:item:d}d) \emph{Under the hypotheses of c), if $f$ is moreover a biregular isomorphism, then it is also an analytic isomorphism.}

By definition, a Z-closed subset of $\mathbf{C}^n$ is defined by the vanishing of a certain number of polynomials. Since a polynomial is continuous for the usual topology (respectively, holomorphic), this proves a) (respectively, b)). To prove c), we may suppose that $\mathrm{U}'=\mathbf{C}$; we must then show that every regular function on $\mathrm{U}$ is holomorphic, which again follows from the fact that a polynomial is a holomorphic function. Finally, d) is an immediate consequence of c) applied to $f^{-1}$.

Let $\mathrm{X}$ now be an algebraic variety over $\mathbf{C}$ (in the sense of \GAGACite{18}, no.~34, and thus not necessarily irreducible). Let $\mathrm{V}$ be a Z-open subset of $\mathrm{X}$ possessing an algebraic chart

\begin{equation*}
\varphi:\quad \mathrm{V}\longrightarrow\mathrm{U}
\end{equation*}

onto a Z-locally closed subset $\mathrm{U}$ of an affine space. By \hyperref[gaga:lemma:1:item:b]{Lemma 1(b)}, $\mathrm{U}$ may be endowed with a structure of analytic space.

\GAGATarget{gaga:prop:2}\textsc{Proposition 2}. --- \emph{There exists a unique structure of analytic space on $\mathrm{X}$ such that, for every chart $\varphi:\mathrm{V}\to\mathrm{U}$, the Z-open subset $\mathrm{V}$ is open and $\varphi$ is an analytic isomorphism from $\mathrm{V}$ (endowed with the analytic structure induced by that of $\mathrm{X}$) onto $\mathrm{U}$ (endowed with the analytic structure defined in \hyperref[gaga:no:1]{no.~1}).}

(More briefly: every algebraic chart must be an analytic chart.)

Uniqueness is clear, since $\mathrm{X}$ can be covered by Z-open subsets $\mathrm{V}$ possessing charts. To prove existence, let $\varphi:\mathrm{V}\to\mathrm{U}$ be a chart, and transport the analytic structure of $\mathrm{U}$ to $\mathrm{V}$ by means of $\varphi^{-1}$. If $\varphi':\mathrm{V}'\to\mathrm{U}'$ is another chart, the analytic structures induced on $\mathrm{V}\cap\mathrm{V}'$ by $\mathrm{V}$ and by $\mathrm{V}'$ are the same, by \hyperref[gaga:lemma:1:item:d]{Lemma 1(d)}; moreover, $\mathrm{V}\cap\mathrm{V}'$ is open both in $\mathrm{V}$ and in $\mathrm{V}'$, by \hyperref[gaga:lemma:1:item:a]{Lemma 1(a)}. By gluing, we thus obtain on $\mathrm{X}$ a topology and a sheaf $\mathcal{H}_{\mathrm{X}}$ that visibly satisfy axiom $(\mathrm{H}_{\mathrm{I}})$. To verify $(\mathrm{H}_{\mathrm{II}})$, we use axiom $(\mathrm{VA}'_{\mathrm{II}})$ of \GAGACite{18}, no.~34. With the notation of that axiom, the graphs $\mathrm{T}_{ij}$ of the identification relations between two $\mathrm{U}_i$ and $\mathrm{U}_j$ are Z-closed in $\mathrm{U}_i\times\mathrm{U}_j$, hence \emph{a fortiori} closed; this says precisely that $\mathrm{X}$ is Hausdorff.

\emph{Remark}. --- The analytic structure of $\mathrm{X}$ may also be defined directly, without using the charts $\varphi:\mathrm{V}\to\mathrm{U}$.

Define the topology to be the coarsest one for which regular functions on Z-open subsets of $\mathrm{X}$ are continuous, and define $\mathcal{H}_{x,\mathrm{X}}$ to be the analytic subring of $\mathcal{C}(\mathrm{X})_x$ generated by $\mathcal{O}_{x,\mathrm{X}}$ (in the sense of \GAGACite{3}, Expos\'e VIII). We leave it to the reader to verify that the two definitions are equivalent.

Henceforth, $\mathrm{X}^h$ will denote the set $\mathrm{X}$ endowed with the structure of analytic space just defined. The topology of $\mathrm{X}^h$ is \emph{finer} than the topology of $\mathrm{X}$; since $\mathrm{X}^h$ can be covered by finitely many open subsets possessing charts, $\mathrm{X}^h$ is a locally compact space that is \emph{countable at infinity}.

The following properties follow immediately from the definition of $\mathrm{X}^h$.

If $\mathrm{X}$ and $\mathrm{Y}$ are two algebraic varieties, then $(\mathrm{X}\times\mathrm{Y})^h=\mathrm{X}^h\times\mathrm{Y}^h$. If $\mathrm{Y}$ is a Z-locally closed subset of $\mathrm{X}$, then $\mathrm{Y}^h$ is an analytic subset of $\mathrm{X}^h$; moreover, the analytic structure of $\mathrm{Y}^h$ agrees with the structure induced on $\mathrm{Y}$ by $\mathrm{X}^h$. Finally, if $f:\mathrm{X}\to\mathrm{Y}$ is a regular map of an algebraic variety $\mathrm{X}$ into an algebraic variety $\mathrm{Y}$, then $f$ is also a holomorphic map from $\mathrm{X}^h$ to $\mathrm{Y}^h$.

\GAGASubsection{6. Relations between the local ring at a point and the ring of holomorphic functions at that point}{gaga:no:6}

Let $\mathrm{X}$ be an algebraic variety, and let $x$ be a point of $\mathrm{X}$. We shall compare the local ring $\mathcal{O}_x$ of regular functions on $\mathrm{X}$ at $x$ with the local ring $\mathcal{H}_x$ of holomorphic functions on $\mathrm{X}^h$ in a neighbourhood of $x$.

Since every regular function is holomorphic, every function $f\in\mathcal{O}_x$ defines a germ of a holomorphic function at $x$, which we denote by $\theta(f)$. The map $\theta:\mathcal{O}_x\to\mathcal{H}_x$ is a homomorphism and maps the maximal ideal of $\mathcal{O}_x$ into that of $\mathcal{H}_x$; it therefore extends continuously to a homomorphism $\hat\theta:\hat{\mathcal{O}}_x\to\hat{\mathcal{H}}_x$ from the completion of $\mathcal{O}_x$ to that of $\mathcal{H}_x$ (cf. \GAGAInternalRef{gaga:appendix}{Appendix}, \GAGAInternalRef{gaga:no:24}{no.~24}).

\GAGATarget{gaga:prop:3}\textsc{Proposition 3}. --- \emph{The homomorphism $\hat\theta:\hat{\mathcal{O}}_x\to\hat{\mathcal{H}}_x$ is bijective.}

We shall prove the preceding proposition together with another result.

Let $\mathrm{Y}$ be a Z-locally closed subset of $\mathrm{X}$, and let $\mathcal{J}_x(\mathrm{Y})$ (or $\mathcal{J}_x(\mathrm{Y},\mathrm{X})$ when the ambient space $\mathrm{X}$ is to be specified) be the ideal of $\mathcal{O}_x$ consisting of functions $f$ whose restriction to $\mathrm{Y}$ vanishes in a Z-neighbourhood of $x$ (cf. \GAGACite{18}, no.~39). The image of $\mathcal{J}_x(\mathrm{Y})$ under $\theta$ is clearly contained in the ideal $\mathcal{A}_x(\mathrm{Y})$ of $\mathcal{H}_x$ defined in \hyperref[gaga:no:3]{no.~3}.

\GAGATarget{gaga:prop:4}\textsc{Proposition 4}. --- \emph{The ideal $\mathcal{A}_x(\mathrm{Y})$ is generated by $\theta(\mathcal{J}_x(\mathrm{Y}))$.}

We shall first prove Propositions \GAGAInternalRef{gaga:prop:3}{3} and \GAGAInternalRef{gaga:prop:4}{4} in the special case where $\mathrm{X}$ is the affine space $\mathbf{C}^n$. Proposition \GAGAInternalRef{gaga:prop:3}{3} is then trivial, since $\hat{\mathcal{O}}_x$ and $\hat{\mathcal{H}}_x$ are both the algebra $\mathbf{C}[[z_1,\ldots,z_n]]$ of formal power series in $n$ indeterminates. Turn now to Proposition \GAGAInternalRef{gaga:prop:4}{4}, and let $\mathfrak{a}$ be the ideal of $\mathcal{H}_x$ generated by $\mathcal{J}_x(\mathrm{Y})$ (the ring $\mathcal{O}_x$ being identified with a subring of $\mathcal{H}_x$ by means of $\theta$). Every ideal of $\mathcal{H}_x$ defines a germ of an analytic subset of $\mathrm{X}$ at $x$ (cf. \GAGACite{1}, no.~3, or \GAGACite{3}, Expos\'e VI, p.~6); it is clear that the germ defined by $\mathfrak{a}$ is precisely $\mathrm{Y}$. Let $f$ be an element of $\mathcal{A}_x(\mathrm{Y})$. By the ``theorem of zeros'' (which is valid for ideals of $\mathcal{H}_x$; cf. \GAGACite{14}, p.~278, as well as \GAGACite{2}, Expos\'e XIV, p.~3, and \GAGACite{3}, Expos\'e VIII, p.~9), there exists an integer $r\geqslant0$ such that $f^r\in\mathfrak{a}$. \emph{A fortiori},

\begin{equation*}
f^r\in\mathfrak{a}\hat{\mathcal{H}}_x
=\mathcal{J}_x(\mathrm{Y})\hat{\mathcal{H}}_x
=\mathcal{J}_x(\mathrm{Y})\hat{\mathcal{O}}_x.
\end{equation*}

But the ideal $\mathcal{J}_x(\mathrm{Y})$ is an intersection of prime ideals corresponding to the irreducible components of $\mathrm{Y}$ passing through $x$. By a theorem of \textsc{Chevalley} (cf. \GAGACite{16}, p.~40, and \GAGACite{17}, p.~67), the same is therefore true of the ideal $\mathcal{J}_x(\mathrm{Y})\hat{\mathcal{O}}_x$, and the relation $f^r\in\mathcal{J}_x(\mathrm{Y})\hat{\mathcal{O}}_x$ consequently implies $f\in\mathcal{J}_x(\mathrm{Y})\hat{\mathcal{O}}_x$. Since $\mathcal{H}_x$ is a Noetherian local ring, $\mathfrak{a}\hat{\mathcal{H}}_x\cap\mathcal{H}_x=\mathfrak{a}$ (cf. \GAGACite{15}, Chapter IV, or \GAGAInternalRef{gaga:appendix}{Appendix}, \GAGAInternalRef{gaga:prop:27}{Proposition 27}); hence $f\in\mathfrak{a}$, which proves Proposition \GAGAInternalRef{gaga:prop:4}{4} in the case under consideration.

Pass to the general case. Since the question is local, we may suppose that $\mathrm{X}$ is a subvariety of an affine space, which we denote by $\mathrm{U}$. By definition,

\begin{equation*}
\mathcal{O}_x=\mathcal{O}_{x,\mathrm{U}}/\mathcal{J}_x(\mathrm{X},\mathrm{U})
\qquad\text{and}\qquad
\mathcal{H}_x=\mathcal{H}_{x,\mathrm{U}}/\mathcal{A}_x(\mathrm{X},\mathrm{U}).
\end{equation*}

The map $\theta:\mathcal{O}_x\to\mathcal{H}_x$ is obtained on passing to the quotient from the map $\theta:\mathcal{O}_{x,\mathrm{U}}\to\mathcal{H}_{x,\mathrm{U}}$. From what precedes, we know that $\hat\theta:\hat{\mathcal{O}}_{x,\mathrm{U}}\to\hat{\mathcal{H}}_{x,\mathrm{U}}$ is bijective and that $\mathcal{A}_x(\mathrm{X},\mathrm{U})=\theta(\mathcal{J}_x(\mathrm{X},\mathrm{U}))\mathcal{H}_{x,\mathrm{U}}$. Proposition \GAGAInternalRef{gaga:prop:3}{3} follows immediately by applying \GAGAInternalRef{gaga:prop:29}{Proposition 29} of the \GAGAInternalRef{gaga:appendix}{Appendix}. As for Proposition \GAGAInternalRef{gaga:prop:4}{4}, it follows because $\mathcal{A}_x(\mathrm{Y})$ is the canonical image of the ideal $\mathcal{A}_x(\mathrm{Y},\mathrm{U})$, which is generated by $\theta(\mathcal{J}_x(\mathrm{Y},\mathrm{U}))$, by what precedes.

Proposition \GAGAInternalRef{gaga:prop:3}{3} shows in particular that $\theta:\mathcal{O}_x\to\mathcal{H}_x$ is \emph{injective}, which allows us to identify $\mathcal{O}_x$ with the subring $\theta(\mathcal{O}_x)$ of $\mathcal{H}_x$. With this identification understood, we have:

\GAGATarget{gaga:cor:1}\textsc{Corollary 1}. --- \emph{The pair of rings $(\mathcal{O}_x,\mathcal{H}_x)$ is a flat pair (in the sense of \GAGAInternalRef{gaga:appendix}{Appendix}, \GAGAInternalRef{gaga:def:4}{Definition 4}).}

This is an immediate consequence of \GAGAInternalRef{gaga:prop:3}{Proposition 3} and \GAGAInternalRef{gaga:prop:28}{Proposition 28} of the \GAGAInternalRef{gaga:appendix}{Appendix}.

\GAGATarget{gaga:cor:2}\textsc{Corollary 2}. --- \emph{The rings $\mathcal{O}_x$ and $\mathcal{H}_x$ have the same dimension.}

Indeed, the dimension of a Noetherian local ring is known to equal that of its completion (cf. \GAGACite{16}, p.~26).

In view of the results stated in \hyperref[gaga:no:4]{no.~4}, we obtain the following result (where we suppose $X$ irreducible in order to simplify the statement):

\GAGATarget{gaga:cor:3}\textsc{Corollary 3}. --- \emph{If $X$ is an irreducible algebraic variety of dimension $r$, the analytic space $X^h$ has analytic dimension $r$ at each of its points.}

\GAGASubsection{7. Relations between the usual topology and the Zariski topology of an algebraic variety}{gaga:no:7}

\GAGATarget{gaga:prop:5}\textsc{Proposition 5}. --- \emph{Let $X$ be an algebraic variety and $U$ a subset of $X$. If $U$ is Z-open and Z-dense in $X$, then $U$ is dense in $X$.}

Let $Y$ be the complement of $U$ in $X$; it is a Z-closed subset of $X$. Let $x$ be a point of $X$. If $x$ were not in the closure of $U$, then $Y=X$ in a neighbourhood of $x$, whence $\mathcal{A}_x(Y)=0$, with the notation of \hyperref[gaga:no:6]{no.~6}. Since $\mathcal{A}_x(Y)$ contains $\theta(\mathcal{J}_x(Y))$ and $\theta$ is injective (\GAGAInternalRef{gaga:prop:3}{Proposition 3}), we would have $\mathcal{J}_x(Y)=0$. This would mean that $Y=X$ in a Z-neighbourhood of $x$, contrary to the hypothesis that $U$ is Z-dense in $X$, qed.

\emph{Remark}. --- One easily sees that \GAGAInternalRef{gaga:prop:5}{Proposition 5} is \emph{equivalent} to the injectivity of $\theta:\mathcal{O}_x\to\mathcal{H}_x$, a much more elementary fact than \GAGAInternalRef{gaga:prop:3}{Proposition 3}, which may, for example, be proved by reduction to the case of a curve.

We now give two simple applications of \GAGAInternalRef{gaga:prop:5}{Proposition 5}.

\GAGATarget{gaga:prop:6}\textsc{Proposition 6}. --- \emph{An algebraic variety $X$ is complete if and only if it is compact.}

First recall a result of \textsc{Chow} (cf. \GAGACite{7} and \GAGACite{19}, no.~4): for every algebraic variety $X$, there exist a projective variety $Y$, a subset $U$ of $Y$ that is Z-open and Z-dense in $Y$, and a surjective regular map $f:U\to X$ whose graph $T$ is Z-closed in $X\times Y$. We have $U=Y$ if and only if $X$ is complete.

Suppose first that $X$ is complete. Then $X=f(Y)$, and since every projective variety is compact for the usual topology, it follows that $X$ is compact. Conversely, suppose that $X$ is compact. Then so is $T$, which is closed in $X\times Y$, and hence so is $U$, since it is the projection of $T$ onto $Y$. Thus $U$ is closed in $Y$, and \GAGAInternalRef{gaga:prop:5}{Proposition 5} shows that $U=Y$, completing the proof.

The following lemma is essentially due to \textsc{Chevalley}.

\GAGATarget{gaga:lemma:2}\textsc{Lemma 2}. --- \emph{Let $f:X\to Y$ be a regular map from an algebraic variety $X$ to an algebraic variety $Y$, and suppose that $f(X)$ is Z-dense in $Y$. Then there exists a subset $U\subset f(X)$ that is Z-open and Z-dense in $Y$.}

When $X$ and $Y$ are irreducible, this result is well known; cf., for example, \GAGACite{4}, Expos\'e 3, or \GAGACite{17}, p.~15. We reduce the general case to that one. Let $X_i$, $i\in I$, be the irreducible components of $X$, and let $Y_i$ be the Z-closure of $f(X_i)$ in $Y$. The $Y_i$ are irreducible and $Y=\bigcup Y_i$; hence there exists $J\subset I$ such that the $Y_j$, $j\in J$, are the irreducible components of $Y$. By the result just recalled, for every $j\in J$ there exists a subset $U_j\subset f(X_j)$ that is Z-open and Z-dense in $Y_j$. After shrinking $U_j$, we may moreover suppose that $U_j$ meets none of the $Y_k$, $k\in J$, $k\neq j$. Then $U=\bigcup_{j\in J}U_j$ is a subset of $Y$ having all the required properties.

\GAGATarget{gaga:prop:7}\textsc{Proposition 7}. --- \emph{If $f:X\to Y$ is a regular map from an algebraic variety $X$ to an algebraic variety $Y$, the closure and the Z-closure of $f(X)$ in $Y$ coincide.}

Let $T$ be the Z-closure of $f(X)$ in $Y$. Applying \GAGAInternalRef{gaga:lemma:2}{Lemma 2} to $f:X\to T$, we see that there exists a subset $U\subset f(X)$ that is Z-open and Z-dense in $T$. By \GAGAInternalRef{gaga:prop:5}{Proposition 5}, $U$ is therefore dense in $T$, and \emph{a fortiori} the same is true of $f(X)$. This shows that $T$ is contained in the closure of $f(X)$; the opposite inclusion is clear, and the proof is complete.

\GAGASubsection{8. An analytic criterion for regularity}{gaga:no:8}

Every regular map is known to be holomorphic. The following proposition (which we shall supplement in \GAGAInternalRef{gaga:no:19}{no.~19}) indicates when the converse is true.

\GAGATarget{gaga:prop:8}\textsc{Proposition 8}. --- \emph{Let $X$ and $Y$ be two algebraic varieties, and let $f:X\to Y$ be a holomorphic map from $X$ to $Y$. If the graph $T$ of $f$ is a Z-locally closed subset (that is, an algebraic subvariety) of $X\times Y$, then $f$ is regular.}

Let $p=\operatorname{pr}_X$ be the canonical projection of $T$ onto the first factor $X$ of $X\times Y$. The map $p$ is regular and bijective, and its inverse is the map $x\mapsto(x,f(x))$, which is holomorphic by hypothesis. Thus $p$ is an analytic isomorphism, and it remains only to show that $p$ is a biregular isomorphism (since $f=\operatorname{pr}_Y\circ p^{-1}$). This follows from the next proposition.

\GAGATarget{gaga:prop:9}\textsc{Proposition 9}. --- \emph{Let $T$ and $X$ be two algebraic varieties, and let $p:T\to X$ be a bijective regular map. If $p$ is an analytic isomorphism from $T$ onto $X$, then it is also a biregular isomorphism.}

We first show that $p$ is a homeomorphism for the Zariski topologies of $T$ and $X$. Let $F$ be a Z-closed subset of $T$. Since $p$ is an analytic isomorphism, it is \emph{a fortiori} a homeomorphism, and $p(F)$ is closed in $X$. Applying \GAGAInternalRef{gaga:prop:7}{Proposition 7} to $p:F\to X$, we conclude that $p(F)$ is Z-closed in $X$, proving the assertion.

It remains to show that $p$ carries the sheaf $\mathcal{O}_X$ of local rings of $X$ to the sheaf $\mathcal{O}_T$ of local rings of $T$. More precisely, if $t$ is a point of $T$ and $x=p(t)$, the map $p$ defines a homomorphism

\begin{equation*}
p^*:\quad\mathcal{O}_{x,X}\longrightarrow\mathcal{O}_{t,T},
\end{equation*}

and we must prove that $p^*$ is bijective.\footnote{The proof that follows was communicated to me by P. \textsc{Samuel}.}

Since $p$ is a Z-homeomorphism, $p^*$ is injective, allowing us to identify $\mathcal{O}_{x,X}$ with a subring of $\mathcal{O}_{t,T}$. To simplify notation, put $A=\mathcal{O}_{x,X}$ and $A'=\mathcal{O}_{t,T}$, so that $A\subset A'$. Likewise, let $B$ (respectively, $B'$) denote the ring $\mathcal{H}_{x,X}$ (respectively, $\mathcal{H}_{t,T}$), and regard $A$ and $A'$ as embedded in $B$ and $B'$, respectively (which is legitimate by \GAGAInternalRef{gaga:prop:3}{Proposition 3}). The hypothesis that $p$ is an analytic isomorphism means that $B=B'$.

Let $X_i$ be the irreducible components of $X$ passing through $x$. Each $X_i$ determines a prime ideal $\mathfrak{p}_i=\mathcal{J}_x(X_i)$ of the ring $A$, and the quotient local ring $A_i=A/\mathfrak{p}_i$ is precisely the local ring of $x$ on $X_i$. Thus the field of fractions $K_i$ of $A_i$ is precisely the field of rational functions on the irreducible variety $X_i$. The ideals $\mathfrak{p}_i$ are clearly the minimal prime ideals of $A$, and $0=\bigcap\mathfrak{p}_i$. The set $S$ of elements of $A$ that belong to none of the $\mathfrak{p}_i$ is multiplicatively closed (it is easily seen to be the set of regular elements of $A$). The total ring of fractions $A_S$ is the direct product of the fields $K_i$ (cf. \GAGAInternalRef{gaga:lemma:3}{Lemma 3} below).

Let $T_i=p^{-1}(X_i)$. Since $p$ is a Z-homeomorphism, the $T_i$ are the irreducible components of $T$ passing through $t$ and define prime ideals $\mathfrak{p}'_i$ of $A'$. Put $A'_i=A'/\mathfrak{p}'_i$, and denote by $K'_i$ the field of fractions of $A'_i$. The total ring of fractions $A'_{S'}$ is the direct product of the $K'_i$. Notice that $\mathfrak{p}'_i\cap A=\mathfrak{p}_i$, whence $A_i\subset A'_i$, $K_i\subset K'_i$, and $A_S\subset A'_{S'}$.

We first show that $K_i=K'_i$, in other words, that $p$ defines a ``birational'' correspondence between $T_i$ and $X_i$. Since $p:T_i\to X_i$ is a Z-homeomorphism, $T_i$ and $X_i$ have the same dimension, and the fields $K_i$ and $K'_i$ have the same transcendence degree over $\mathbf{C}$. Put $n_i=[K'_i:K_i]$. It is known\footnote{This is a classical and easily proved result about correspondences. A slightly weaker result, sufficient for the present application, may be found in \GAGACite{17}, p.~16.} that there exists a nonempty Z-open subset $U_i$ of $X_i$ such that the inverse image of every point of $U_i$ consists of exactly $n_i$ points of $T_i$. Since $p$ is bijective, this shows that $n_i=1$, and hence $K_i=K'_i$.

Since $A_S$ (respectively, $A'_{S'}$) is the direct product of the $K_i$ (respectively, the $K'_i$), it follows that $A_S=A'_{S'}$. Let $f'\in A'$. By what precedes, $f'\in A_S$; in other words, there exist $g\in A$ and $s\in S$ such that $g=sf'$. Thus $g\in sA'$, whence $g\in sB'$, that is, $g\in sB$. But by \GAGAInternalRef{gaga:cor:1}{Corollary 1} to \GAGAInternalRef{gaga:prop:3}{Proposition 3}, $(A,B)$ is a flat pair, and therefore $sB\cap A=sA$; cf. \GAGAInternalRef{gaga:appendix}{Appendix}, \GAGAInternalRef{gaga:no:22}{no.~22}. Hence $g\in sA$, so there exists $f\in A$ such that $g=sf$, or equivalently $s(f-f')=0$. Since $s$ is not a zero-divisor in $A'$, this implies $f=f'$, that is, $A=A'$, qed.

We used in the course of the proof the following result, which we now prove.

\GAGATarget{gaga:lemma:3}\textsc{Lemma 3}. --- \emph{Let $A$ be a commutative ring in which the zero ideal is the intersection of finitely many distinct minimal prime ideals $\mathfrak{p}_i$. Let $K_i$ be the field of fractions of $A/\mathfrak{p}_i$, and let $S$ be the set of elements of $A$ that belong to none of the $\mathfrak{p}_i$. Then the ring of fractions $A_S$ is isomorphic to the direct product of the $K_i$.}

The prime ideals of $A_S$ are known to correspond bijectively to the prime ideals of $A$ that do not meet $S$ (cf. \GAGACite{15}, Chapter IV, \S 3, to which we refer for everything concerning rings of fractions). Consequently, if $\mathfrak{m}_i=\mathfrak{p}_iA_S$, then the $\mathfrak{m}_i$ are the only prime ideals of $A_S$; in particular, they are maximal ideals. They are clearly distinct, since $\mathfrak{m}_i\cap A=\mathfrak{p}_i$ (\GAGACite{15}, \emph{loc. cit.}). Moreover, the field $A_S/\mathfrak{m}_i$ is generated by $A/\mathfrak{p}_i$ and therefore agrees with $K_i$. It remains to show that the canonical homomorphism

\begin{equation*}
\varphi:A_S\longrightarrow\prod A_S/\mathfrak{m}_i=\prod K_i
\end{equation*}

is bijective.

First, $\bigcap\mathfrak{p}_i=0$ implies $\bigcap\mathfrak{m}_i=0$, which shows that $\varphi$ is injective. Let $\mathfrak{b}_i$ be the product, in the ring $A_S$, of the ideals $\mathfrak{m}_j$ with $j\neq i$, and put $\mathfrak{b}=\sum\mathfrak{b}_i$. The ideal $\mathfrak{b}$ is contained in none of the $\mathfrak{m}_i$, hence equals $A_S$, and there exist elements $x_i\in\mathfrak{b}_i$ such that $\sum x_i=1$. We have

\begin{equation*}
x_i\equiv1\pmod{\mathfrak{m}_i}
\qquad\text{and}\qquad
x_i\equiv0\pmod{\mathfrak{m}_j},\quad j\neq i,
\end{equation*}

which shows that $\varphi(A_S)$ contains the elements $(1,0,\ldots,0),\ldots,(0,\ldots,0,1)$ of $\prod K_i$. Since these elements generate the $A_S$-module $\prod K_i$, this proves that $\varphi$ is bijective and completes the proof.
